% ===========================================================================
%  Chapitre 7 — La fonction logarithme népérien
%  PE 2026, composante ANALYSE
% ===========================================================================
\bmtchapitre{La fonction logarithme népérien}
  {Connaître la fonction $\ln$, son ensemble de définition et ses propriétés
   algébriques ; résoudre les équations et inéquations du type
   $\ln u(x) = a$, $\ln u(x) = \ln v(x)$ et
   $a\ln^{2}x+b\ln x+c = 0$ ; étudier une fonction contenant un logarithme.}
  {Analyse \textbullet\ environ 5 heures}

Voici la fonction qui donne au problème du baccalauréat sa couleur : sur les
quinze sessions dépouillées pour ce livre, dix reposent sur un logarithme.
Ce n'est pas un hasard. Le logarithme est la fonction qui transforme les
produits en sommes, et c'est cette propriété — écrite en une ligne — qui a
fait sa fortune bien avant les calculatrices.

Son autre intérêt est plus terre à terre : elle offre, en une seule fonction,
tout ce que l'analyse de l'année sait faire. Un ensemble de définition à
justifier, une limite infinie en $0$, une croissance lente en $+\infty$, une
dérivée simple, un tableau de variation, une courbe à tracer.

% ---------------------------------------------------------------------------
\begin{revision}
Uniquement les chapitres 1, 4, 5 et 6.

\begin{enumerate}[label=\textbf{\arabic*.}, leftmargin=*, itemsep=0.35em]
  \item Résoudre $x^{2}-3x+2 = 0$.
  \item Résoudre l'inéquation $2x-6 > 0$.
  \item Dériver $x \mapsto \dfrac{1}{x}$ et $x \mapsto x^{2}-x$.
  \item Pour quelles valeurs de $x$ a-t-on $x-3 > 0$ ?
  \item Calculer $\displaystyle\lim_{x \to 0^{+}} \dfrac{1}{x}$.
  \item Donner la dérivée de $\dfrac1v$ en fonction de $v$ et de $v'$.
\end{enumerate}
\end{revision}

\begin{automatismes}
Après le cours, sans calculatrice.

\begin{enumerate}[label=\textbf{\alph*)}, leftmargin=*, itemsep=0.25em]
  \item $\ln 1$
  \item $\ln e$
  \item $\ln e^{3}$
  \item $\ln \dfrac1e$
  \item $\ln 6-\ln 2$
  \item $\ln 4+\ln 5$
  \item $2\ln 3$ sous la forme $\ln a$
  \item L'ensemble de définition de $x \mapsto \ln(x-5)$
\end{enumerate}
\end{automatismes}

\begin{corrige}[automatismes]
a) $0$ \quad b) $1$ \quad c) $3$ \quad d) $-1$ \quad e) $\ln3$ \quad
f) $\ln20$ \quad g) $\ln9$ \quad h) $\intervalleo{5}{+\infty}$.
\end{corrige}

% ===========================================================================
\section{Définition et premières propriétés}

\begin{definition}[fonction logarithme népérien]
On appelle \textbf{fonction logarithme népérien}, notée $\ln$, l'unique
fonction définie sur $\intervalleo{0}{+\infty}$, dérivable, telle que
\[
  \ln' (x)= \frac1x \quad\text{pour tout } x>0, \qquad \ln 1 = 0 .
\]
\end{definition}

\begin{remarque}
Deux conséquences immédiates, et essentielles. D'abord,
\textbf{$\ln x$ n'existe que pour $x$ strictement positif} : écrire
$\ln(-3)$ ou $\ln 0$ n'a aucun sens. Ensuite, $\ln'(x) = \frac1x > 0$ : la
fonction $\ln$ est \textbf{strictement croissante} sur
$\intervalleo{0}{+\infty}$.
\end{remarque}

\begin{definition}[le nombre $e$]
Puisque $\ln$ est strictement croissante et prend toutes les valeurs réelles,
il existe un unique réel dont le logarithme vaut $1$. On le note $e$ :
\[
  \ln e = 1, \qquad e \approx 2{,}718 .
\]
\end{definition}

\begin{propriete}[limites de $\ln$]
\[
  \lim_{x \to 0^{+}} \ln x = -\infty, \qquad
  \limpinf \ln x = +\infty, \qquad
  \limpinf \frac{\ln x}{x} = 0 .
\]
\end{propriete}

\begin{avertissement}
La troisième limite est celle qu'on oublie. Elle dit que le logarithme
\emph{croît moins vite que $x$} : à l'infini, $x$ « écrase » $\ln x$. Les
sujets la donnent presque toujours en indication — « on admet que
$\limpinf \frac{\ln x}{x} = 0$ » — et attendent qu'on s'en serve pour lever
une indétermination.
\end{avertissement}

\begin{visualisation}[la courbe du logarithme]
\begin{center}
\begin{tikzpicture}
\begin{axis}[
  width = 11cm, height = 6.2cm,
  axis lines = middle, xlabel = {$x$}, ylabel = {$y$},
  xmin = -0.6, xmax = 8.4, ymin = -3.4, ymax = 2.9,
  xtick = {1,2,3,4,5,6,7,8}, ytick = {-3,-2,-1,1,2},
  axis line style = {bmtGris},
  tick label style = {font=\scriptsize, color=bmtGris},
  grid = both, grid style = {bmtGrisClair, line width=0.3pt},
  clip = true, restrict y to domain = -6:4,
]
  \addplot[bmtIndigo, line width=1.5pt, domain=0.05:8.3, samples=250]
    {ln(x)};
  \addplot[bmtLaterite, line width=1pt, dash pattern=on 3pt off 2pt,
           domain=0.05:3.2] {x-1};
  \filldraw[bmtVetiver] (axis cs:1,0) circle (2pt);
  \filldraw[bmtVetiver] (axis cs:2.718,1) circle (2pt);
  \node[bmtVetiver, font=\scriptsize, anchor=north west]
    at (axis cs:2.78,1.02) {$(e\,;\,1)$};
  \node[bmtVetiver, font=\scriptsize, anchor=north west]
    at (axis cs:1.05,-0.05) {$(1\,;\,0)$};
  \node[bmtLaterite, font=\scriptsize, anchor=south east]
    at (axis cs:3.15,2.2) {$y = x-1$};
\end{axis}
\end{tikzpicture}
\end{center}

Trois choses à retenir de cette courbe. Elle n'existe qu'à droite de l'axe
des ordonnées, qui est son asymptote verticale. Elle traverse l'axe des
abscisses en $1$, et passe par le point $(e\,;\,1)$. Enfin, elle monte
toujours, mais de plus en plus lentement : entre $1$ et $2$ elle gagne
$0{,}69$ ; entre $100$ et $101$, elle ne gagne plus que $0{,}01$. La
tangente en $1$, d'équation $y = x-1$, reste au-dessus de la courbe partout.
\end{visualisation}

% ===========================================================================
\section{Les propriétés algébriques}

\begin{theoreme}[la propriété fondamentale]
Pour tous réels strictement positifs $a$ et $b$ :
\[
  \ln(a\times b) = \ln a+\ln b .
\]
\end{theoreme}

\begin{propriete}[les conséquences]
Pour $a>0$, $b>0$ et $n$ entier :
\[
  \ln\frac ab = \ln a-\ln b, \qquad
  \ln\frac1b = -\ln b, \qquad
  \ln\left(a^{n}\right) = n\ln a, \qquad
  \ln\sqrt a = \frac12\ln a .
\]
\end{propriete}

\begin{demonstration}[de $\ln\frac1b = -\ln b$]
En appliquant la propriété fondamentale au produit $b\times\frac1b = 1$ :
\[
  \ln b+\ln\frac1b = \ln 1 = 0,
\]
d'où le résultat. Les autres s'en déduisent de la même façon.
\end{demonstration}

\begin{erreurclassique}
$\ln(a+b)$ n'est \textbf{pas} $\ln a+\ln b$. La propriété transforme un
\emph{produit} en somme, jamais une somme en somme. De même,
$\ln(a-b) \neq \ln a-\ln b$, et $\dfrac{\ln a}{\ln b} \neq \ln\dfrac ab$.
Trois erreurs, la même origine : on applique une formule sans regarder
l'opération sur laquelle elle porte.
\end{erreurclassique}

\begin{exemple}[trois simplifications]
\textbf{a)} $\ln 12 = \ln(4\times3) = \ln 4+\ln 3 = 2\ln 2+\ln 3$.

\textbf{b)} $\ln\dfrac{e^{3}}{5} = 3-\ln 5$.

\textbf{c)} $\ln 8-3\ln 2 = 3\ln 2-3\ln 2 = 0$.
\end{exemple}

\begin{entrainement}[calcul algébrique]
\exo[\niveau{1}] Écrire $\ln 18$ en fonction de $\ln 2$ et $\ln 3$.

\exo[\niveau{1}] Simplifier $\ln 25-2\ln 5$.

\exo[\niveau{2}] Simplifier $\ln\left(e^{2}\right)+\ln\dfrac1e$.

\exo[\niveau{2}] Écrire $\dfrac12\ln 9+\ln 2$ sous la forme $\ln a$.

\exo[\niveau{2}] Montrer que $\ln\dfrac{a^{2}b}{c} = 2\ln a+\ln b-\ln c$
pour $a$, $b$, $c$ strictement positifs.

\exo[\niveau{3}] Sachant que $\ln 2 \approx 0{,}69$ et
$\ln 3 \approx 1{,}10$, donner une valeur approchée de $\ln 24$ et de
$\ln \frac32$.
\end{entrainement}

\begin{corrige}[calcul algébrique]
\textbf{1.} $\ln18 = \ln(2\times9) = \ln2+2\ln3$.
\textbf{2.} $2\ln5-2\ln5 = 0$.
\textbf{3.} $2-1 = 1$.
\textbf{4.} $\frac12\ln9 = \ln3$, donc le total vaut $\ln6$.
\textbf{5.} $\ln\frac{a^{2}b}{c} = \ln\left(a^{2}b\right)-\ln c
= \ln a^{2}+\ln b-\ln c = 2\ln a+\ln b-\ln c$.
\textbf{6.} $\ln24 = 3\ln2+\ln3 \approx 2{,}07+1{,}10 = 3{,}17$ ;
$\ln\frac32 = \ln3-\ln2 \approx 0{,}41$.
\end{corrige}

% ===========================================================================
\section{Équations et inéquations avec un logarithme}

\begin{avertissement}
Toute résolution commence par les \textbf{conditions d'existence} : chaque
expression écrite sous un $\ln$ doit être strictement positive. On les pose
avant de calculer, et l'on rejette à la fin les solutions qui ne les
respectent pas. Une équation résolue sans conditions d'existence perd
systématiquement des points.
\end{avertissement}

\begin{theoreme}[les trois équations du programme]
Pour $u(x)>0$ et $v(x)>0$ :
\[
  \ln u(x) = a \daccord u(x) = e^{a}, \qquad
  \ln u(x) = \ln v(x) \daccord u(x) = v(x).
\]
Et l'équation $a\ln^{2}x+b\ln x+c = 0$ se résout en posant $X = \ln x$ : on
obtient un trinôme du second degré en $X$.
\end{theoreme}

\begin{propriete}[inéquations]
$\ln$ étant strictement croissante, pour $u(x)>0$ et $v(x)>0$ :
\[
  \ln u(x) < \ln v(x) \daccord u(x) < v(x).
\]
Le sens de l'inégalité est donc conservé — c'est ce qui distingue le
logarithme d'une fonction décroissante.
\end{propriete}

\begin{methode}{résoudre une équation en $\ln$}
\begin{enumerate}[label=\textbf{\arabic*.}, leftmargin=*, itemsep=0.15em]
  \item Écrire les conditions d'existence et l'ensemble $\mathcal D$ où l'on
        travaille.
  \item Regrouper les logarithmes d'un même côté, en un seul $\ln$, grâce
        aux propriétés algébriques.
  \item Appliquer $\ln u = \ln v \Rightarrow u = v$, ou passer à
        l'exponentielle si le second membre est un nombre.
  \item Résoudre l'équation obtenue.
  \item \emph{Confronter} les solutions à $\mathcal D$ et conclure.
\end{enumerate}
\end{methode}

\begin{exemple}[une solution rejetée]
Résolvons $\ln(x-1)+\ln(x+2) = \ln 4$.

\emph{Conditions.} $x-1>0$ et $x+2>0$, donc $x>1$ :
$\mathcal D = \intervalleo{1}{+\infty}$.

\emph{Regroupement.} $\ln\left[(x-1)(x+2)\right] = \ln 4$, d'où
$(x-1)(x+2) = 4$, soit $x^{2}+x-6 = 0$.

\emph{Résolution.} $\Delta = 25$, racines $-3$ et $2$.

\emph{Confrontation.} $-3 \notin \mathcal D$ : on la rejette. Reste
$\mathcal S = \{2\}$.
\end{exemple}

\begin{exemple}[un changement d'inconnue]
Résolvons $\ln^{2}x-3\ln x+2 = 0$ sur $\intervalleo{0}{+\infty}$.

En posant $X = \ln x$ : $X^{2}-3X+2 = 0$, de racines $1$ et $2$. Donc
$\ln x = 1$ ou $\ln x = 2$, c'est-à-dire $x = e$ ou $x = e^{2}$.
\end{exemple}

\begin{entrainement}[équations et inéquations]
\exo[\niveau{1}] $\ln x = 3$

\exo[\niveau{1}] $\ln(2x-1) = 0$

\exo[\niveau{2}] $\ln(x+3) = \ln(2x-1)$

\exo[\niveau{2}] $\ln x+\ln(x-3) = \ln 4$

\exo[\niveau{2}] $\ln^{2}x-\ln x = 0$

\exo[\niveau{3}] $2\ln^{2}x+5\ln x-3 = 0$

\exo[\niveau{3}] $\ln(x-1) \leq \ln(5-x)$

\exo[\niveau{3}] $\ln x < 2$
\end{entrainement}

\begin{corrige}[équations et inéquations]
\textbf{1.} $x = e^{3}$.
\textbf{2.} $2x-1 = 1$ donc $x = 1$ (condition $x>\frac12$ respectée).
\textbf{3.} Conditions : $x>\frac12$. Alors $x+3 = 2x-1$ donne $x = 4$.
\textbf{4.} Conditions : $x>3$. Alors $x(x-3) = 4$, soit $x^{2}-3x-4 = 0$,
de racines $-1$ et $4$ ; on rejette $-1$ : $\mathcal S = \{4\}$.
\textbf{5.} $\ln x(\ln x-1) = 0$ : $x = 1$ ou $x = e$.
\textbf{6.} $X = \ln x$ : $2X^{2}+5X-3 = 0$, $\Delta = 49$, $X = -3$ ou
$X = \frac12$. D'où $x = e^{-3}$ ou $x = \sqrt e$.
\textbf{7.} Conditions : $1<x<5$. L'inégalité équivaut à $x-1 \leq 5-x$,
soit $x \leq 3$ : $\mathcal S = \intervalleof{1}{3}$.
\textbf{8.} Condition $x>0$ ; $\ln x<\ln e^{2}$ donne $x<e^{2}$ :
$\mathcal S = \intervalleo{0}{e^{2}}$.
\end{corrige}

\begin{qcm}
\exo $\ln(3x)$ est défini pour :
\textbf{a)} $x \geq 0$ \quad \textbf{b)} $x>0$ \quad \textbf{c)} $x \neq 0$
\quad \textbf{d)} tout réel \reponseqcm{b}

\exo $\ln 2+\ln 3$ vaut :
\textbf{a)} $\ln 5$ \quad \textbf{b)} $\ln 6$ \quad \textbf{c)} $\ln 23$
\quad \textbf{d)} $6$ \reponseqcm{b}

\exo L'équation $\ln x = -2$ a pour solution :
\textbf{a)} aucune \quad \textbf{b)} $x = e^{2}$ \quad
\textbf{c)} $x = e^{-2}$ \quad \textbf{d)} $x = -e^{2}$ \reponseqcm{c}

\exo $\displaystyle\lim_{x \to 0^{+}} \ln x$ vaut :
\textbf{a)} $0$ \quad \textbf{b)} $1$ \quad \textbf{c)} $+\infty$ \quad
\textbf{d)} $-\infty$ \reponseqcm{d}

\exo La dérivée de $x \mapsto \ln(3x+1)$ sur
$\intervalleo{-\frac13}{+\infty}$ est :
\textbf{a)} $\dfrac{1}{3x+1}$ \quad \textbf{b)} $\dfrac{3}{3x+1}$ \quad
\textbf{c)} $\dfrac{1}{3}$ \quad \textbf{d)} $3\ln(3x+1)$ \reponseqcm{b}
\end{qcm}

% ===========================================================================
\section{Étudier une fonction contenant un logarithme}

\begin{propriete}[dérivée]
Si $u$ est dérivable et strictement positive sur un intervalle $I$, alors
$x \mapsto \ln u(x)$ est dérivable sur $I$ et
\[
  \left(\ln u\right)'(x) = \frac{u'(x)}{u(x)} .
\]
\end{propriete}

\begin{exemple}[l'étude type du baccalauréat]
Soit $f(x) = x-2\ln x$ sur $\intervalleo{0}{+\infty}$.

\emph{Limites.} En $0^{+}$ : $x \to 0$ et $-2\ln x \to +\infty$, donc
$f(x) \to +\infty$ ; la droite $x = 0$ est asymptote verticale. En
$+\infty$, on écrit
\[
  f(x) = x\left(1-\frac{2\ln x}{x}\right),
\]
et comme $\frac{\ln x}{x} \to 0$, la parenthèse tend vers $1$ :
$f(x) \to +\infty$.

\emph{Dérivée.} $f'(x) = 1-\dfrac{2}{x} = \dfrac{x-2}{x}$. Sur
$\intervalleo{0}{+\infty}$, le dénominateur est positif : $f'$ a le signe de
$x-2$.

\emph{Variations.} $f$ décroît sur $\intervalleof{0}{2}$, croît sur
$\intervallefo{2}{+\infty}$, et son minimum vaut
$f(2) = 2-2\ln 2 \approx 0{,}61$.

Ce minimum étant strictement positif, la courbe reste entièrement au-dessus
de l'axe des abscisses : l'équation $f(x) = 0$ n'a aucune solution.
\end{exemple}

\begin{methode}{lever l'indétermination $\infty-\infty$ avec un logarithme}
On met $x$ en facteur et l'on fait apparaître $\dfrac{\ln x}{x}$, dont la
limite en $+\infty$ est nulle. C'est le geste attendu par tous les sujets :
\[
  x-2\ln x = x\left(1-2\,\frac{\ln x}{x}\right),
  \qquad
  \frac{\ln x}{x}\xrightarrow[x\to+\infty]{}0 .
\]
\end{methode}

\begin{entrainement}[fonctions avec logarithme]
\exo[\niveau{1}] Dériver $f(x) = \ln(2x+5)$ sur
$\intervalleo{-\frac52}{+\infty}$.

\exo[\niveau{2}] Dériver $g(x) = x\ln x$ sur $\intervalleo{0}{+\infty}$ et
étudier le signe de $g'$.

\exo[\niveau{2}] Dériver $h(x) = \dfrac{\ln x}{x}$ et déterminer en quel
point la courbe admet une tangente horizontale.

\exo[\niveau{2}] $f(x) = \ln\left(x^{2}+1\right)$ : donner $\Df$, la dérivée
et le sens de variation.

\exo[\niveau{3}] $f(x) = 2x-1-2\ln x$ sur $\intervalleo{0}{+\infty}$ :
calculer $f'$, dresser le tableau de variation et en déduire le signe de
$f$.
\end{entrainement}

\begin{corrige}[fonctions avec logarithme]
\textbf{1.} $\frac{2}{2x+5}$.
\textbf{2.} $g'(x) = \ln x+1$, qui s'annule en $e^{-1}$ : négative avant,
positive après.
\textbf{3.} $h'(x) = \frac{1-\ln x}{x^{2}}$, nulle pour $x = e$ : la
tangente est horizontale au point $\left(e\,;\,\frac1e\right)$.
\textbf{4.} $\Df = \R$ ; $f'(x) = \frac{2x}{x^{2}+1}$, du signe de $x$ : la
fonction décroît sur $\R^{-}$, croît sur $\R^{+}$, minimum $f(0) = 0$.
\textbf{5.} $f'(x) = 2-\frac2x = \frac{2(x-1)}{x}$ : négative sur
$\intervalleo{0}{1}$, positive après. Le minimum vaut $f(1) = 1 > 0$, donc
$f(x)>0$ pour tout $x>0$.
\end{corrige}

% ===========================================================================
\begin{annales}
Le fonds est ici abondant : le problème du baccalauréat de la série~A repose
sur un logarithme à la majorité des sessions. Les énoncés suivants sont
repris tels quels, à l'exception des questions faisant appel à la primitive
ou à l'aire, traitées au chapitre 9.

\exoann{Baccalauréat, Madagascar, série A, session 2016 — problème}
Soit $f(x) = x-\ln x$, de courbe $(C)$ dans un repère orthonormé d'unité
$1$ cm.
\begin{enumerate}[label=\textbf{\alph*)}, leftmargin=*, itemsep=0.15em]
  \item Déterminer l'ensemble de définition $\Df$ de $f$.
  \item Vérifier que, pour tout $x>0$,
        $f(x) = x\left(1-\dfrac{\ln x}{x}\right)$ ; en admettant que
        $\limpinf \dfrac{\ln x}{x} = 0$, calculer $\limpinf f(x)$.
  \item Calculer $\displaystyle\lim_{x \to 0^{+}} f(x)$ et en déduire une
        asymptote.
  \item Démontrer que, pour tout $x>0$, $f'(x) = \dfrac{x-1}{x}$, puis
        dresser le tableau de variation de $f$.
  \item Compléter le tableau de valeurs pour $x = 2$ ; $e$ ; $4$ ; $6$, à
        $10^{-2}$ près, puis tracer $(C)$.
\end{enumerate}

\exoann{Baccalauréat, Madagascar, série A, session 2012 — problème}
Soit $f(x) = 2x-1-2\ln x$, de courbe $(C)$ dans un repère orthonormé
d'unité $2$ cm.
\begin{enumerate}[label=\textbf{\alph*)}, leftmargin=*, itemsep=0.15em]
  \item Vérifier que $\displaystyle\lim_{x \to 0^{+}} f(x) = +\infty$.
  \item Montrer que, pour tout $x>0$,
        $f(x) = x\left(2-\dfrac1x-2\,\dfrac{\ln x}{x}\right)$, puis en
        déduire $\limpinf f(x)$.
  \item Calculer $f'(x)$ et dresser le tableau de variation.
\end{enumerate}

\exoann{Baccalauréat, Madagascar, série A, session 2010 — problème}
Soit $f(x) = x-2\ln x$ sur $\intervalleo{0}{+\infty}$, de courbe $(C)$ dans
un repère orthonormé d'unité $2$ cm.
\begin{enumerate}[label=\textbf{\alph*)}, leftmargin=*, itemsep=0.15em]
  \item Vérifier que $f(x) = x\left(1-2\,\dfrac{\ln x}{x}\right)$ pour tout
        $x>0$.
  \item En admettant que $\limpinf \dfrac{\ln x}{x} = 0$, déterminer
        $\limpinf f(x)$, puis étudier la branche infinie de la courbe.
  \item Calculer $f'(x)$, étudier son signe, dresser le tableau de
        variation.
\end{enumerate}

\exoann{Baccalauréat, Madagascar, série A, session 2004 — problème}
Soit $f(x) = \dfrac1x-1+\ln x$ sur $\intervalleo{0}{+\infty}$.
\begin{enumerate}[label=\textbf{\alph*)}, leftmargin=*, itemsep=0.15em]
  \item Calculer $f'(x)$ et étudier son signe suivant les valeurs de $x$.
  \item Calculer $\limpinf f(x)$.
  \item Montrer que $f(x) = \dfrac{1-x+x\ln x}{x}$, puis, en admettant que
        $\displaystyle\lim_{x \to 0^{+}} x\ln x = 0$, en déduire
        $\displaystyle\lim_{x \to 0^{+}} f(x)$.
  \item Dresser le tableau de variation de $f$.
\end{enumerate}

\exoann{Baccalauréat, Madagascar, série A, session 2015 — problème}
Soit $f(x) = \left(\ln x\right)^{2}-1$ sur $\intervalleo{0}{+\infty}$.
\begin{enumerate}[label=\textbf{\alph*)}, leftmargin=*, itemsep=0.15em]
  \item Calculer la limite de $f$ en $0^{+}$ et l'interpréter
        graphiquement.
  \item Calculer la limite de $f$ en $+\infty$.
  \item Montrer que $f'(x) = \dfrac{2\ln x}{x}$, étudier son signe et
        dresser le tableau de variation.
  \item Vérifier que $f(x) = (\ln x-1)(\ln x+1)$, résoudre $f(x) = 0$, et en
        déduire les coordonnées des points d'intersection de la courbe avec
        l'axe des abscisses.
\end{enumerate}

\exoann{Baccalauréat, Madagascar, série A, session 2020 — problème}
Soit $f(x) = \ln(1+x)-\ln(3-x)$ sur $\intervalleo{-1}{3}$.
\begin{enumerate}[label=\textbf{\alph*)}, leftmargin=*, itemsep=0.15em]
  \item Prouver que $\displaystyle\lim_{x \to (-1)^{+}} f(x) = -\infty$ et
        donner une interprétation graphique.
  \item Vérifier que $\displaystyle\lim_{x \to 3^{-}} f(x) = +\infty$ ;
        qu'en conclut-on pour la courbe ?
  \item Montrer que, pour tout $x$ de $\intervalleo{-1}{3}$,
        $f'(x) = \dfrac{4}{(1+x)(3-x)}$, puis dresser le tableau de
        variation.
  \item Déterminer les coordonnées du point d'intersection de la courbe avec
        l'axe des abscisses, puis écrire l'équation de la tangente au point
        d'abscisse $x_{0} = 1$.
\end{enumerate}

\exoann{Baccalauréat, Madagascar, série A, session 2002 — problème
(extrait)}
Soit $f(x) = 2\ln x\,(\ln x-1)$.
\begin{enumerate}[label=\textbf{\alph*)}, leftmargin=*, itemsep=0.15em]
  \item Déterminer $\Df$ et les limites aux bornes de $\Df$.
  \item Montrer que $f'(x) = \dfrac{2}{x}\left(2\ln x-1\right)$, puis
        dresser le tableau de variation.
  \item Résoudre $f(x) = 0$.
\end{enumerate}

\exoann{Baccalauréat, Madagascar, série A, session 2000 — problème
(extrait)}
Soit $f(x) = \ln(x+4)-\ln(2-x)$ sur $\intervalleo{-4}{2}$.
\begin{enumerate}[label=\textbf{\alph*)}, leftmargin=*, itemsep=0.15em]
  \item Calculer les limites de $f$ en $-4$ et en $2$ ; interpréter
        graphiquement.
  \item Montrer que $f'(x) = \dfrac{6}{(x+4)(2-x)}$ et dresser le tableau de
        variation.
  \item Déterminer le point d'intersection de la courbe avec l'axe des
        abscisses, puis écrire l'équation de la tangente en ce point.
\end{enumerate}
\end{annales}

\begin{corrige}[annales]
\textbf{1.} $\Df = \intervalleo{0}{+\infty}$. La factorisation donne
$f(x) = x\left(1-\frac{\ln x}{x}\right) \to +\infty$ puisque la parenthèse
tend vers $1$. En $0^{+}$, $-\ln x \to +\infty$ donc $f \to +\infty$ :
l'axe des ordonnées est asymptote verticale. $f'(x) = 1-\frac1x
= \frac{x-1}{x}$, négative sur $\intervalleo{0}{1}$, positive après :
minimum $f(1) = 1$. Valeurs : $f(2) \approx 1{,}31$ ; $f(e) \approx 1{,}72$ ;
$f(4) \approx 2{,}61$ ; $f(6) \approx 4{,}21$.

\textbf{2.} En $0^{+}$, $-2\ln x \to +\infty$ et $2x-1 \to -1$ : la limite
est $+\infty$. La factorisation donne une parenthèse de limite $2$, donc
$f \to +\infty$. $f'(x) = 2-\frac2x = \frac{2(x-1)}{x}$ : minimum
$f(1) = 1$.

\textbf{3.} $f(x) = x\left(1-2\frac{\ln x}{x}\right) \to +\infty$ ; comme
$\frac{f(x)}{x} \to 1$ et $f(x)-x = -2\ln x \to -\infty$, la courbe admet
une branche infinie de direction asymptotique la droite $y = x$, sans
asymptote. $f'(x) = \frac{x-2}{x}$ : minimum $f(2) = 2-2\ln2 \approx 0{,}61$.

\textbf{4.} $f'(x) = -\frac{1}{x^{2}}+\frac1x = \frac{x-1}{x^{2}}$ :
négative sur $\intervalleo{0}{1}$, positive après. En $+\infty$,
$\frac1x \to 0$ et $\ln x \to +\infty$ : $f \to +\infty$. La factorisation
$f(x) = \frac{1-x+x\ln x}{x}$ donne, en $0^{+}$, un numérateur qui tend vers
$1$ et un dénominateur vers $0^{+}$ : $f \to +\infty$. Minimum
$f(1) = 0$ — la courbe touche l'axe des abscisses en $x = 1$ sans le
traverser.

\textbf{5.} En $0^{+}$, $\ln x \to -\infty$ donc $(\ln x)^{2} \to +\infty$ et
$f \to +\infty$ : asymptote verticale $x = 0$. En $+\infty$,
$f \to +\infty$ également. $f'(x) = 2\ln x\times\frac1x
= \frac{2\ln x}{x}$, du signe de $\ln x$ : $f$ décroît sur
$\intervalleo{0}{1}$, croît après ; minimum $f(1) = -1$. Enfin
$(\ln x)^{2}-1 = (\ln x-1)(\ln x+1)$ s'annule pour $\ln x = \pm1$, soit
$x = e$ et $x = e^{-1}$ : les points d'intersection sont
$A\left(\frac1e\,;\,0\right)$ et $B(e\,;\,0)$.

\textbf{6.} En $(-1)^{+}$, $\ln(1+x) \to -\infty$ et $\ln(3-x) \to \ln4$ :
$f \to -\infty$, asymptote verticale $x = -1$. En $3^{-}$,
$-\ln(3-x) \to +\infty$ : $f \to +\infty$, asymptote verticale $x = 3$. La
dérivée vaut
$\frac{1}{1+x}+\frac{1}{3-x} = \frac{(3-x)+(1+x)}{(1+x)(3-x)}
= \frac{4}{(1+x)(3-x)}$, strictement positive sur l'intervalle : $f$ est
strictement croissante. $f(x) = 0$ équivaut à $1+x = 3-x$, soit $x = 1$ ; le
point est $(1\,;\,0)$. Comme $f'(1) = \frac{4}{2\times2} = 1$, la tangente a
pour équation $y = x-1$.

\textbf{7.} $\Df = \intervalleo{0}{+\infty}$ ; en $0^{+}$, $\ln x \to
-\infty$ et $\ln x-1 \to -\infty$ : le produit tend vers $+\infty$. En
$+\infty$, les deux facteurs tendent vers $+\infty$ : $f \to +\infty$. En
développant, $f(x) = 2(\ln x)^{2}-2\ln x$, donc
$f'(x) = \frac{4\ln x}{x}-\frac2x = \frac2x\left(2\ln x-1\right)$, qui
s'annule pour $\ln x = \frac12$, soit $x = \sqrt e$ ; minimum
$f(\sqrt e) = 2\times\frac12\times\left(-\frac12\right) = -\frac12$. Enfin
$f(x) = 0$ pour $\ln x = 0$ ou $\ln x = 1$, soit $x = 1$ ou $x = e$.

\textbf{8.} En $-4^{+}$, $\ln(x+4) \to -\infty$ : $f \to -\infty$ ; en
$2^{-}$, $-\ln(2-x) \to +\infty$ : $f \to +\infty$. Les droites $x = -4$ et
$x = 2$ sont asymptotes verticales.
$f'(x) = \frac{1}{x+4}+\frac{1}{2-x} = \frac{(2-x)+(x+4)}{(x+4)(2-x)}
= \frac{6}{(x+4)(2-x)} > 0$ sur $\intervalleo{-4}{2}$ : $f$ y est
strictement croissante. $f(x) = 0$ donne $x+4 = 2-x$, soit $x = -1$ ; comme
$f'(-1) = \frac{6}{3\times3} = \frac23$, la tangente en $I(-1\,;\,0)$ a pour
équation $y = \frac23(x+1)$.
\end{corrige}

% ===========================================================================
\begin{vraifaux}
\exo $\ln(x^{2})$ et $2\ln x$ ont le même ensemble de définition.

\exo Si $\ln a = \ln b$ alors $a = b$.

\exo $\ln$ est croissante, donc $\ln x > 0$ pour tout $x$ de son ensemble de
définition.

\exo L'équation $\ln x = -5$ n'a pas de solution.

\exo $\ln(2x) = 2\ln x$ pour tout $x>0$.
\end{vraifaux}

\begin{corrige}[vrai ou faux]
\textbf{1.} Faux — $\ln(x^{2})$ est défini pour $x \neq 0$, donc aussi pour
$x<0$ ; $2\ln x$ seulement pour $x>0$.
\textbf{2.} Vrai — la stricte croissance assure l'injectivité.
\textbf{3.} Faux — $\ln x<0$ pour $0<x<1$.
\textbf{4.} Faux — la solution est $x = e^{-5}$, un nombre positif très
petit.
\textbf{5.} Faux — $\ln(2x) = \ln2+\ln x$, alors que $2\ln x = \ln(x^{2})$.
\end{corrige}

% ===========================================================================
\begin{jecherche}[le carbone et le temps]
La quantité de carbone 14 restant dans un échantillon après $t$ milliers
d'années est $Q(t) = Q_{0}\,\times\,0{,}886^{\,t}$, où $Q_{0}$ est la
quantité initiale.

\begin{enumerate}[label=\textbf{\arabic*.}, leftmargin=*, itemsep=0.3em]
  \item Quelle proportion de carbone reste-t-il après mille ans ? après
        cinq mille ans ?
  \item On cherche la \emph{demi-vie} : le temps $t$ au bout duquel la
        moitié du carbone a disparu. Écrire l'équation correspondante.
  \item Montrer que $t = \dfrac{\ln 0{,}5}{\ln 0{,}886}$, puis en donner une
        valeur approchée. On donne $\ln 0{,}5 \approx -0{,}69$ et
        $\ln 0{,}886 \approx -0{,}121$.
  \item Un os retrouvé contient $30\,\%$ du carbone 14 d'un os vivant.
        Estimer son âge.
\end{enumerate}
\end{jecherche}

\begin{corrige}[le carbone et le temps]
\textbf{1.} Après mille ans, $88{,}6\,\%$ ; après cinq mille ans,
$0{,}886^{5} \approx 0{,}546$, soit environ $54{,}6\,\%$.

\textbf{2.} On cherche $t$ tel que $Q_{0}\times0{,}886^{t}
= \frac{Q_{0}}{2}$, soit $0{,}886^{t} = 0{,}5$.

\textbf{3.} En prenant le logarithme des deux membres :
$t\ln 0{,}886 = \ln 0{,}5$, d'où
$t = \frac{\ln0{,}5}{\ln0{,}886} \approx \frac{-0{,}69}{-0{,}121}
\approx 5{,}7$. La demi-vie est d'environ $5\,700$ ans.

\textbf{4.} $0{,}886^{t} = 0{,}3$ donne
$t = \frac{\ln0{,}3}{\ln0{,}886} \approx \frac{-1{,}20}{-0{,}121}
\approx 9{,}9$ : l'os a environ dix mille ans.
\end{corrige}

% ===========================================================================
\begin{jemeteste}
Trente minutes.

\begin{enumerate}[label=\textbf{\arabic*.}, leftmargin=*, itemsep=0.3em]
  \item Simplifier $\ln 27-3\ln 3$ et $\ln\dfrac{e^{4}}{e}$.
  \item Résoudre $\ln(2x+1) = \ln(x+4)$.
  \item Résoudre $\ln^{2}x+\ln x-6 = 0$.
  \item Donner l'ensemble de définition de $x \mapsto \ln(4-x^{2})$.
  \item Dériver $f(x) = 3x-\ln x$ et dresser son tableau de variation sur
        $\intervalleo{0}{+\infty}$.
\end{enumerate}
\end{jemeteste}

\begin{corrige}[je me teste]
\textbf{1.} $3\ln3-3\ln3 = 0$ ; et $4-1 = 3$.
\textbf{2.} Conditions $x>-\frac12$ ; $2x+1 = x+4$ donne $x = 3$.
\textbf{3.} $X^{2}+X-6 = 0$ : $X = -3$ ou $X = 2$, d'où $x = e^{-3}$ ou
$x = e^{2}$.
\textbf{4.} $4-x^{2}>0$ donne $\Df = \intervalleo{-2}{2}$.
\textbf{5.} $f'(x) = 3-\frac1x = \frac{3x-1}{x}$ : négative sur
$\intervalleo{0}{\frac13}$, positive après ; minimum
$f\left(\frac13\right) = 1+\ln3 \approx 2{,}1$.
\end{corrige}

% ===========================================================================
\begin{commeaubac}
\bmtsource{Baccalauréat, Madagascar, série A, session 2026 — problème
\textbullet\ 10 points}
On considère la fonction numérique $f$ définie sur
$\intervalleo{0}{+\infty}$ par
\[
  f(x) = -\frac1x-2x+\ln x .
\]
On désigne par $(\mathcal C)$ la courbe représentative de $f$ dans un repère
orthonormé $\left(O\,;\,\veci\,,\,\vecj\right)$ d'unité graphique $1$ cm.

\begin{enumerate}[label=\textbf{\arabic*.}, leftmargin=*, itemsep=0.25em]
  \item Calculer les limites de $f$ aux bornes de son ensemble de
        définition. \bareme{1}
  \item \begin{enumerate}[label=\textbf{\alph*)}, leftmargin=1.4em, itemsep=0.15em]
      \item Pour tout $x \in \intervalleo{0}{+\infty}$, calculer $f'(x)$.
            \bareme{1,5}
      \item Montrer que $f'(x) = \dfrac{-2x^{2}+x+1}{x^{2}}$ pour tout
            $x>0$. \bareme{0,75}
      \item En déduire le signe de $f'(x)$ pour tout
            $x \in \intervalleo{0}{+\infty}$. \bareme{0,5}
    \end{enumerate}
  \item Dresser le tableau de variation de $f$. \bareme{1,5}
  \item Écrire une équation de la tangente $(T)$ à $(\mathcal C)$ au point
        d'abscisse $x_{0} = 1$. \bareme{0,75}
  \item Tracer $(\mathcal C)$ et $(T)$ sur l'intervalle
        $\intervalleof{0}{4}$. \bareme{1,5}
\end{enumerate}

\begin{remarque}
La question 6 de ce sujet porte sur une primitive et sur un calcul d'aire :
elle figure au chapitre 9, une fois l'intégrale introduite.
\end{remarque}
\end{commeaubac}

\begin{corrige}[comme au baccalauréat — Madagascar A 2026]
\textbf{1.} En $0^{+}$ : $-\frac1x \to -\infty$, $-2x \to 0$ et
$\ln x \to -\infty$, donc $f(x) \to -\infty$ ; l'axe des ordonnées est
asymptote verticale. En $+\infty$ : $-\frac1x \to 0$, et
$-2x+\ln x = x\left(-2+\frac{\ln x}{x}\right) \to -\infty$ puisque la
parenthèse tend vers $-2$. Donc $f(x) \to -\infty$.

\textbf{2. a)} $f'(x) = \frac{1}{x^{2}}-2+\frac1x$.
\textbf{b)} En réduisant au même dénominateur $x^{2}$ :
\[
  f'(x) = \frac{1-2x^{2}+x}{x^{2}} = \frac{-2x^{2}+x+1}{x^{2}} .
\]
\textbf{c)} Le dénominateur est strictement positif. Le trinôme
$-2x^{2}+x+1$ a pour discriminant $1+8 = 9$ et pour racines
$\frac{-1-3}{-4} = 1$ et $\frac{-1+3}{-4} = -\frac12$ ; son coefficient
dominant étant négatif, il est positif entre les racines. Sur
$\intervalleo{0}{+\infty}$ : $f'(x)>0$ sur $\intervalleo{0}{1}$,
$f'(1) = 0$, et $f'(x)<0$ sur $\intervalleo{1}{+\infty}$.

\textbf{3.} $f$ croît sur $\intervalleof{0}{1}$, décroît sur
$\intervallefo{1}{+\infty}$. Le maximum vaut
$f(1) = -1-2+0 = -3$, atteint en $x = 1$ ; les deux limites valent
$-\infty$.

\textbf{4.} $f(1) = -3$ et $f'(1) = 0$ : la tangente est horizontale,
$(T) : y = -3$.

\textbf{5.} On place l'asymptote verticale (l'axe des ordonnées), la
tangente horizontale $y = -3$, le sommet $(1\,;\,-3)$, puis quelques
points : $f(0{,}5) \approx -3{,}69$, $f(2) \approx -3{,}81$,
$f(4) \approx -6{,}86$. La courbe monte jusqu'au sommet puis redescend, en
restant toujours sous la droite $y = -3$.
\end{corrige}

\begin{notepourlaclasse}
Cinq heures, et pas une de plus, alors que le chapitre porte à lui seul la
moitié des points du problème. La seule façon d'y arriver est de renoncer à
construire le logarithme : on le \emph{pose} par sa dérivée, comme le fait le
programme, et l'on passe aussitôt aux propriétés algébriques.

Deux automatismes à installer coûte que coûte, parce qu'ils reviennent à
chaque session : la factorisation par $x$ qui fait apparaître
$\frac{\ln x}{x}$, et le réflexe des conditions d'existence avant toute
résolution d'équation.

Le bloc d'annales est volontairement long — huit sujets réels. Il n'est pas
fait pour être traité en entier : choisissez-en trois, dont 2026 et 2020, et
laissez les autres en révision autonome. Un élève qui aura vu passer six de
ces huit énoncés reconnaîtra la structure du problème le jour de l'épreuve.
\end{notepourlaclasse}
